\newpage\SourceLeaf{48}\ArithmeticHead{41}{R}{DIE ANALYSIS RESIDUORUM.}\HistoricalFoot{X\textsubscript{2}.}{6}
\begin{center}\bfseries II.\quad Gauss’ arithmetische Abhandlungen und Nachlass.\end{center}
\begin{center}\bfseries a.\quad\textls[100]{Die Analysis Residuorum.}\end{center}
\begin{center}\bfseries 16.\end{center}
\noindent\hspace*{\parindent}An die D. A. fügen wir nun naturgemäss die Betrachtungen der A. R., deren weitere Ausführung \textsc{Gauss} als Fortsetzung der D. A. geplant hatte. Im ersten Teile der A. R. (W. II, S. 199—211) wird die Lösung der binomischen Kongruenz \(x^n\equiv1\) in Bezug auf einen Primzahlmodul \(p\) gelehrt. Man darf \(n\) als Teiler von \(p-1\) und als eine Primzahlpotenz voraussetzen; ist \(\varphi(n)=a^\alpha b^\beta\ldots\), so kommt die Auflösung der Kongruenz auf \(\alpha\) Kongruenzen vom Grade \(a\), \(\beta\) Kongruenzen vom Grade \(b,\ldots\) zurück. Dies zeigt sich für eine Primzahl \(n\) in ganz entsprechender Weise wie bei der Gleichung \(x^n=1\). Sei \(r\) primitive Wurzel von
\[
x^n\equiv1\ (\mathrm{mod.}\ p),
\]
\(\rho\) primitive Wurzel von
\[
x^{\varphi(n)}\equiv1\ (\mathrm{mod.}\ n)
\]
und \(\varphi(n)=e\mathbin{.}f\), und bildet man die Perioden
\[
\begin{gathered}
\eta_i=r^\rho+r^{\rho^{e+i}}+r^{\rho^{2e+i}}+\ldots+r^{\rho^{(f-1)e+i}},\\
(i=0,1,2,\ldots,e-1)
\end{gathered}
\]
so ist
\[
\eta_i\eta_k\equiv C+m_0\eta_0+\ldots+m_{e-1}\eta_{e-1}\ (\mathrm{mod.}\ p),
\]
wo \(C\) und die \(m\) ganze Zahlen bedeuten. Setzt man ferner \(e=e'f'\) und zerlegt \(\eta_i\) in kleinere Perioden
\[
\eta_i=\eta'_{0i}+\eta'_{1i}+\ldots+\eta'_{e'-1,i},
\]
so ist die Summe
\[
\eta'_{0h}\mathbin{.}\eta'_{0k}+\eta'_{1h}\mathbin{.}\eta'_{1k}+\ldots+\eta'_{e'-1,h}\mathbin{.}\eta'_{e'-1,k}
\]
eine lineare Funktion von den \(\eta_i\). Daher kann man, wenn die letzteren bekannt sind, die Potenzsummen
\newpage\SourceLeaf{49}\ArithmeticHead{42}{L}{BACHMANN, ÜBER GAUSS’ ZAHLENTHEORETISCHE ARBEITEN.}
\[
\eta_{0i}'{}^m+\eta_{1i}'{}^m+\ldots\eta_{e-1,i}'{}^m
\]
bilden und mittels ihrer eine Kongruenz, deren Wurzeln die kleineren Perioden sind. Die Fortsetzung dieses Verfahrens führt das obgenannte Ergebnis herbei. Dabei ist zu beachten, dass aus einer Periode \(\eta_0\) die übrigen mittels linearer Kongruenzen von der Gestalt
\[
\begin{gathered}
\eta_0^h\equiv C_h+m_0^{(h)}\eta_0+m_1^{(h)}\eta_1+\ldots+m_{e-1}^{(h)}\eta_{e-1}\\
(h=0,1,2,\ldots,e-1)
\end{gathered}
\]
bestimmt werden können, vorausgesetzt, dass die Determinante dieser Kongruenzen nicht durch \(p\) teilbar ist. Die hier bleibende Lücke hat \textsc{Gauss}, wenn seine Tagebuchnotiz vom 21. Juli 1797 (Nr. 68 T.) hierauf zu beziehen ist, zu ergänzen gewusst, indem er die Theorie der Kongruenzen zu Hilfe zog, die in Bezug auf einen Primzahlpotenzmodul gedacht werden; in welcher Weise aber, das tritt nicht zu Tage.

Von der besonderen Kongruenz \(x^n\equiv1\) wendet sich \textsc{Gauss} (W. II, S. 212—242) allgemein zur Betrachtung der Kongruenzen höheren Grades
\[
F(x)\equiv0\ (\mathrm{mod.}\ p).
\]
Da die Aufgabe, eine Wurzel einer solchen zu finden, nur ein spezieller Fall der allgemeinen ist, die Funktion \(F(x)\ (\mathrm{mod.}\ p)\) in ihre einfachsten Faktoren zu zerlegen, so sieht man sich zu Untersuchungen geführt, welche denjenigen der Theorie der rationalen ganzen Zahlen völlig analog sind, indem die ganzen Funktionen \(F(x)\) an Stelle der letztern treten. In solcher Analogie ist später die Theorie systematisch von \textsc{Dedekind}\footnote{\textsc{Crelles} Journal für Mathematik \textbf{54} (1857), S. 1.}\ entwickelt worden, nachdem schon früher \textsc{Schönemann}\begingroup\renewcommand{\thefootnote}{**)}\footnote{\textsc{Crelles} Journal für Mathematik \textbf{31} (1846), S. 269; \textbf{32} (1846), S. 93.}\endgroup\ sie auf anderen Grundlagen aufgebaut hatte.

\textsc{Gauss}’ Darstellung derselben nimmt mehr \textsc{Dedekinds} Weg und liefert bereits die grösste Anzahl der allgemeinen Sätze, zu denen Dieser gelangt ist. Nach Definition der Teilbarkeit einer Funktion \(F(x)\ (\mathrm{mod.}\ p)\) treten die Primfunktionen als Grundelemente hervor, in die jede andere Funktion eindeutig zerlegbar ist; für relative Primfunktionen \(A(x),B(x)\) lassen sich andere Funk-
\newpage\SourceLeaf{50}\ArithmeticHead{43}{R}{DIE ANALYSIS RESIDUORUM.}
\noindent tionen \(C(x),D(x)\) finden derart, dass
\[
A(x)\mathbin{.}C(x)+B(x)\mathbin{.}D(x)\equiv1\ (\mathrm{mod.}\ p)
\]
(s. die Notiz vom 19. Aug. 1796, Nr. 27 T.). Die Anzahl der inkongruenten Funktionen \(m^{\mathrm{ten}}\) Grades \((\mathrm{mod.}\ p)\) ist \(p^m\). Die Anzahl \((m)\) der inkongruenten Primfunktionen \(m^{\mathrm{ten}}\) Grades hat \textsc{Gauss} (am 26. Aug. 1797, Nr. 75 T.) sehr einfach auf analytischem Wege mittels einer sie erzeugenden Funktion bestimmt, nachdem er sie früher auf zwei andern Wegen, deren einer in der Handschrift der A. R. mitgeteilt ist, umständlicher erhalten hatte, und hat so die Formel gefunden
\[
p^m=\textstyle\sum d\mathbin{.}(d),
\]
in welcher die Summe auf alle Teiler \(d\) von \(m\) sich erstreckt und aus deren Umkehrung \((m)\) sich ergibt, eine Formel, welche, wenn \(m\) als Primzahl gedacht wird, einen neuen Beweis des \textsc{Fermat}schen Satzes herbeiführt.
